\chapter{Bẫy Của Việc Làm Hài Lòng}
\markboth{Bẫy Của Việc Làm Hài Lòng}{The Trap of Pleasing Everyone}

\section{Gật Đầu Cho Khỏi Mất Lòng}
\begin{truyen}{Gật Đầu Cho Khỏi Mất Lòng}{Saying Yes Just to Avoid Discomfort}
\chuhoa{C}{ó người trẻ rất sợ không khí trở nên khó chịu nên hễ ai nhờ gì cũng gật đầu.}
\textit{(Some young people fear tension so much that they say yes to almost every request.)}

Ban đầu họ thấy mình tử tế và dễ thương trong mắt mọi người.
\textit{(At first, they feel kind and pleasant in the eyes of others.)}

Nhưng càng nhận nhiều phần việc, càng hứa nhiều điều, họ càng sống trong cảm giác bị kéo đi khắp nơi.
\textit{(But the more tasks and promises they accept, the more they feel dragged in every direction.)}

Làm hài lòng người khác bằng cách bỏ mặc giới hạn của mình thường không dẫn tới bình an, mà dẫn tới kiệt sức.
\textit{(Pleasing others by abandoning your own limits rarely brings peace. It usually brings exhaustion.)}

\end{truyen}

\begin{baihoc}
Không phải cái gật đầu nào cũng là lòng tốt. Có cái chỉ là nỗi sợ bị ghét.
\textit{(Not every yes is kindness. Some are simply the fear of being disliked.)}
\end{baihoc}

\begin{ghinhoanh}
\textbf{Short English to remember:}

Not every yes is kindness.

Some are simply the fear of being disliked.

\end{ghinhoanh}

\begin{tuhocnhanh}
\textbf{Câu hỏi tự học / Self-study questions}

1. Nhân vật đang giúp thật hay đang né cảm giác khó xử? \textit{Is the character truly helping or avoiding discomfort?}

2. Cái giá của việc gật đầu liên tục là gì? \textit{What is the cost of constant agreement?}

3. Em cần nói không ở việc nào? \textit{Where do you need to say no?}
\end{tuhocnhanh}
\ngancach

\section{Cố Sống Đúng Ý Gia Đình}
\begin{truyen}{Cố Sống Đúng Ý Gia Đình}{Trying to Live Exactly as the Family Wants}
\chuhoa{N}{hiều người trẻ yêu gia đình thật, nhưng dần dần lại sống như đang thi để được chấp thuận.}
\textit{(Many young people genuinely love their family, yet slowly begin living like they are taking a test for approval.)}

Họ chọn ngành, chọn việc, chọn nhịp sống chỉ để bớt nghe thất vọng từ người thân.
\textit{(They choose majors, jobs, and even lifestyles mainly to avoid hearing disappointment from their relatives.)}

Điều đau là họ có thể làm đúng ý nhà trong một thời gian, nhưng lại ngày càng xa mình.
\textit{(The painful part is that they may satisfy the family for a while while drifting further away from themselves.)}

Hiếu thảo không có nghĩa là giao luôn tay lái cuộc đời cho người khác.
\textit{(Respecting your family does not mean handing over the steering wheel of your life.)}

\end{truyen}

\begin{baihoc}
Yêu gia đình là thật. Nhưng sống thay cho kỳ vọng của người khác quá lâu sẽ làm mình cạn dần.
\textit{(Loving your family is real. But living too long inside someone else’s expectations slowly empties you.)}
\end{baihoc}

\begin{ghinhoanh}
\textbf{Short English to remember:}

Loving your family is real.

But living too long inside someone else’s expectations slowly empties you.

\end{ghinhoanh}

\begin{tuhocnhanh}
\textbf{Câu hỏi tự học / Self-study questions}

1. Nhân vật đang sống bằng giá trị hay bằng sự chấp thuận? \textit{Is the character living by values or by approval?}

2. Điều gì đang bị mất dần? \textit{What is being slowly lost?}

3. Em có cần nói rõ lựa chọn của mình hơn không? \textit{Do you need to explain your own choice more clearly?}
\end{tuhocnhanh}
\ngancach

\section{Tử Tế Đến Mức Tự Xóa Mình}
\begin{truyen}{Tử Tế Đến Mức Tự Xóa Mình}{Being So Nice That You Erase Yourself}
\chuhoa{C}{ó người luôn nhường lời, nhường phần, nhường thời gian, và thấy ngại nếu phải giữ lấy điều mình cần.}
\textit{(Some people always give up their turn, their share, and their time, feeling guilty whenever they protect what they need.)}

Sự tử tế của họ làm ai cũng dễ chịu, trừ chính họ.
\textit{(Their kindness makes everyone comfortable except themselves.)}

Lâu ngày, họ không còn phân biệt đâu là cho đi đẹp, đâu là tự làm mình mờ đi.
\textit{(Over time, they stop seeing the difference between generous giving and self-erasure.)}

Một người không biết giữ phần chính đáng của mình rất dễ bị quen tay bỏ qua.
\textit{(A person who never protects their rightful space is easy for others to keep overlooking.)}

\end{truyen}

\begin{baihoc}
Ranh giới không làm em bớt tử tế. Nó giúp lòng tốt của em bền hơn.
\textit{(Boundaries do not make you less kind. They make your kindness more sustainable.)}
\end{baihoc}

\begin{ghinhoanh}
\textbf{Short English to remember:}

Boundaries do not make you less kind.

They make your kindness more sustainable.

\end{ghinhoanh}

\begin{tuhocnhanh}
\textbf{Câu hỏi tự học / Self-study questions}

1. Nhân vật đang cho đi hay đang tự xóa mình? \textit{Is the character giving or erasing themselves?}

2. Ranh giới giúp điều gì? \textit{What do boundaries protect?}

3. Em đang mệt vì tử tế hay vì không có giới hạn? \textit{Are you tired from kindness or from lacking limits?}
\end{tuhocnhanh}
\ngancach

\section{Làm Đẹp Lòng Mọi Người, Bỏ Lỡ Việc Quan Trọng}
\begin{truyen}{Làm Đẹp Lòng Mọi Người, Bỏ Lỡ Việc Quan Trọng}{Keeping Everyone Happy While Missing What Matters}
\chuhoa{C}{ó bạn trẻ mỗi ngày đều bận trả lời, dàn xếp, giữ hòa khí, và cố để không ai phật ý.}
\textit{(Some young people spend each day replying, smoothing things over, keeping peace, and trying not to upset anyone.)}

Họ tưởng đó là trưởng thành và khéo sống.
\textit{(They think this is maturity and social skill.)}

Nhưng việc quan trọng nhất của đời mình lại cứ bị dời xuống sau cùng: học cho chắc, làm cho xong, nghỉ cho hồi lại.
\textit{(But the most important work of their own life keeps getting pushed to the end: learning well, finishing real work, and recovering properly.)}

Người luôn bận làm vừa lòng đám đông thường không còn đủ sức để xây một đời sống có trục.
\textit{(People who stay busy pleasing the crowd often lose the strength to build a centered life.)}

\end{truyen}

\begin{baihoc}
Nếu mọi người đều ổn mà đời em cứ trượt đi, cách sống đó đang có vấn đề.
\textit{(If everyone is satisfied while your own life keeps slipping, that way of living has a problem.)}
\end{baihoc}

\begin{ghinhoanh}
\textbf{Short English to remember:}

If everyone is satisfied while your own life keeps slipping, that way of living has a problem.

\end{ghinhoanh}

\begin{tuhocnhanh}
\textbf{Câu hỏi tự học / Self-study questions}

1. Nhân vật đang giữ hòa khí hay đang đánh rơi trọng tâm? \textit{Is the character keeping peace or losing focus?}

2. Việc quan trọng nào đang bị dời lại? \textit{What important work is being delayed?}

3. Em cần ưu tiên lại điều gì? \textit{What do you need to prioritize again?}
\end{tuhocnhanh}
\ngancach

\section{Xin Lỗi Cả Khi Mình Không Sai}
\begin{truyen}{Xin Lỗi Cả Khi Mình Không Sai}{Apologizing Even When You Are Not Wrong}
\chuhoa{C}{ó người trẻ quen miệng xin lỗi mỗi khi không khí hơi nặng xuống, dù lỗi không thực sự thuộc về mình.}
\textit{(Some young people apologize automatically whenever the atmosphere gets tense, even when the fault is not truly theirs.)}

Ban đầu điều đó giúp mọi chuyện trôi qua nhanh và giữ được vẻ yên ổn.
\textit{(At first, that helps things pass quickly and preserves a sense of peace.)}

Nhưng lâu dần, họ tập cho người khác quen với việc cứ đẩy áp lực sang mình là xong.
\textit{(But over time, they train others to believe that pressure can always be shifted onto them.)}

Xin lỗi để giữ hòa khí có thể đẹp trong một tình huống, nhưng lặp lại quá nhiều sẽ làm sự thật và lòng tự trọng cùng mờ đi.
\textit{(Apologizing to keep peace may be graceful in one situation, but repeated too often it makes both truth and self-respect fade.)}

\end{truyen}

\begin{baihoc}
Xin lỗi đúng lúc là trưởng thành. Xin lỗi thay cho mọi người là tự làm mình nhỏ lại.
\textit{(Apologizing at the right time is maturity. Apologizing in everyone else's place is self-erasure.)}
\end{baihoc}

\begin{ghinhoanh}
\textbf{Short English to remember:}

Apologizing in everyone else's place is self-erasure.

\end{ghinhoanh}

\begin{tuhocnhanh}
\textbf{Câu hỏi tự học / Self-study questions}

1. Nhân vật đang xin lỗi vì điều gì? \textit{What is the character apologizing for?}

2. Sự lặp lại này làm mờ điều gì? \textit{What does this repetition begin to blur?}

3. Em có đang xin lỗi quá dễ để giữ yên không? \textit{Do you apologize too quickly just to keep peace?}
\end{tuhocnhanh}
\ngancach

\section{Nhận Vai Người Cứu Trong Mọi Câu Chuyện}
\begin{truyen}{Nhận Vai Người Cứu Trong Mọi Câu Chuyện}{Trying to Be the Rescuer in Every Story}
\chuhoa{C}{ó người thấy ai buồn cũng muốn gánh, ai rối cũng muốn đứng ra cứu, ai ngã cũng muốn là người kéo dậy.}
\textit{(Some people feel the need to carry anyone who is sad, rescue anyone in confusion, and lift anyone who falls.)}

Họ tin rằng như vậy mới là sống có tình và có ích.
\textit{(They believe that this is what it means to live with love and usefulness.)}

Nhưng không phải câu chuyện nào mình cũng có quyền bước vào giữa, và không phải ai cũng thật sự muốn được cứu theo cách đó.
\textit{(But not every story is yours to enter, and not everyone truly wants to be rescued in that way.)}

Khi việc giúp người trở thành cách để thấy mình có giá trị, lòng tốt rất dễ biến thành một kiểu lệ thuộc tinh vi.
\textit{(When helping others becomes a way to feel valuable, kindness can turn into a subtle form of dependency.)}

\end{truyen}

\begin{baihoc}
Không phải việc tốt nào cũng cần em đứng giữa. Có lúc lùi lại mới là giúp đúng.
\textit{(Not every good act requires you to stand in the middle. Sometimes stepping back is the truer help.)}
\end{baihoc}

\begin{ghinhoanh}
\textbf{Short English to remember:}

Sometimes stepping back is the truer help.

\end{ghinhoanh}

\begin{tuhocnhanh}
\textbf{Câu hỏi tự học / Self-study questions}

1. Nhân vật đang giúp hay đang cần cảm giác mình quan trọng? \textit{Is the character helping, or needing to feel important?}

2. Vì sao không phải chuyện nào cũng nên chen vào? \textit{Why should you not step into every story?}

3. Em có đang quen làm người cứu hộ quá mức không? \textit{Are you too used to playing the rescuer?}
\end{tuhocnhanh}
\ngancach

\section{Sợ Nói Thật Nên Cứ Nói Cho Vừa Tai}
\begin{truyen}{Sợ Nói Thật Nên Cứ Nói Cho Vừa Tai}{Avoiding Truth and Saying What People Like}
\chuhoa{N}{hiều người trẻ biết rõ điều cần nói, nhưng lại chọn câu dễ nghe hơn vì sợ làm ai đó khó chịu.}
\textit{(Many young people know what needs to be said, but choose the softer version because they fear making someone uncomfortable.)}

Câu nói vừa tai giúp buổi hôm đó yên ổn hơn.
\textit{(The pleasing sentence makes that day calmer.)}

Nhưng nếu kéo dài, sự thật bị dồn lại, còn vấn đề thật thì cứ âm thầm lớn lên.
\textit{(But if this continues, truth gets postponed while the real problem quietly grows.)}

Người quen nói điều vừa lòng thường mất dần khả năng nói điều đúng lúc cần thiết nhất.
\textit{(Those who get used to saying what pleases others gradually lose the ability to say what matters most.)}

\end{truyen}

\begin{baihoc}
Nói thật với sự tử tế khó hơn nói vừa tai, nhưng nó cứu được nhiều điều hơn.
\textit{(Speaking truth with kindness is harder than being pleasant, but it saves more.)}
\end{baihoc}

\begin{ghinhoanh}
\textbf{Short English to remember:}

Truth with kindness saves more than comfort with avoidance.

\end{ghinhoanh}

\begin{tuhocnhanh}
\textbf{Câu hỏi tự học / Self-study questions}

1. Nhân vật đang tránh câu nói nào? \textit{What sentence is the character avoiding?}

2. Điều gì lớn lên khi sự thật bị dời lại? \textit{What grows when truth is delayed?}

3. Em cần nói điều gì một cách tử tế nhưng rõ ràng? \textit{What do you need to say kindly but clearly?}
\end{tuhocnhanh}
\ngancach

\section{Giữ Một Mối Quan Hệ Bằng Cách Tự Ép Mình}
\begin{truyen}{Giữ Một Mối Quan Hệ Bằng Cách Tự Ép Mình}{Forcing Yourself Just to Keep a Relationship}
\chuhoa{C}{ó người ở lại trong một tình bạn, một nhóm, hay một mối quen chỉ vì sợ rời ra sẽ bị xem là thay đổi.}
\textit{(Some people stay in a friendship, a group, or a connection because they fear being judged if they pull away.)}

Để giữ mối quan hệ đó, họ ép mình cười, ép mình đi cùng, ép mình gật đầu với những điều không còn hợp.
\textit{(To keep that relationship, they force themselves to smile, join in, and agree with things that no longer fit.)}

Sự miễn cưỡng kéo dài khiến lòng mình mòn đi trước khi mối quan hệ thật sự tan.
\textit{(Long-term reluctance wears down the heart before the relationship actually ends.)}

Giữ một mối quan hệ bằng cách phản bội cảm giác thật của mình thường chỉ kéo dài hình thức, không giữ được chất lượng.
\textit{(Keeping a relationship by betraying your real feelings often preserves the form, not the quality.)}

\end{truyen}

\begin{baihoc}
Mối quan hệ đáng giữ không bắt em phải liên tục ép mình.
\textit{(A relationship worth keeping does not require you to keep forcing yourself.)}
\end{baihoc}

\begin{ghinhoanh}
\textbf{Short English to remember:}

A relationship worth keeping does not require constant self-forcing.

\end{ghinhoanh}

\begin{tuhocnhanh}
\textbf{Câu hỏi tự học / Self-study questions}

1. Nhân vật đang ép mình để giữ điều gì? \textit{What is the character forcing themselves to preserve?}

2. Điều gì bị mòn đi trước? \textit{What wears down first?}

3. Em có đang giữ một quan hệ chỉ còn hình thức không? \textit{Are you keeping a relationship that is only form now?}
\end{tuhocnhanh}
\ngancach

\section{Chờ Được Thích Rồi Mới Dám Sống Thật}
\begin{truyen}{Chờ Được Thích Rồi Mới Dám Sống Thật}{Waiting for Approval Before Living Honestly}
\chuhoa{N}{hiều người không dám bộc lộ lựa chọn thật, sở thích thật, hay quan điểm thật vì sợ bị mất cảm tình.}
\textit{(Many people do not dare show their real choices, interests, or opinions because they fear losing approval.)}

Họ nghĩ cứ được yêu mến trước đã, rồi sau này sống thật cũng chưa muộn.
\textit{(They think they should be liked first, and only later, perhaps, live honestly.)}

Nhưng càng kéo dài, hình ảnh người khác quen với em lại càng xa con người thật của em.
\textit{(But the longer this lasts, the more the version people know drifts away from who you really are.)}

Đợi đủ được thích rồi mới sống thật thường dẫn tới một ngày không biết phần thật của mình đang ở đâu nữa.
\textit{(Waiting to be approved before living honestly often leads to a day when you can no longer find your real self.)}

\end{truyen}

\begin{baihoc}
Được yêu mến bằng một phiên bản không phải mình là một kiểu cô đơn rất kín.
\textit{(Being liked for a version that is not really you is a quiet kind of loneliness.)}
\end{baihoc}

\begin{ghinhoanh}
\textbf{Short English to remember:}

Being liked for a false version of yourself is a quiet loneliness.

\end{ghinhoanh}

\begin{tuhocnhanh}
\textbf{Câu hỏi tự học / Self-study questions}

1. Nhân vật đang giấu phần thật nào? \textit{What real part is the character hiding?}

2. Vì sao sự được thích này lại cô đơn? \textit{Why can this approval still feel lonely?}

3. Em muốn người khác quý phiên bản nào của mình? \textit{Which version of yourself do you want others to know?}
\end{tuhocnhanh}
\ngancach

\section{Đo Giá Trị Mình Bằng Mức Được Thích}
\begin{truyen}{Đo Giá Trị Mình Bằng Mức Được Thích}{Measuring Your Worth by How Much You Are Liked}
\chuhoa{C}{ó người bước vào đâu cũng âm thầm kiểm xem mình có được thích không, có bị lạnh đi không, có bị bỏ quên không.}
\textit{(Some people quietly measure every room by whether they are liked, ignored, or kept at a distance.)}

Tâm trí họ dần sống bằng phản ứng của người khác nhiều hơn bằng giá trị mình đang theo đuổi.
\textit{(Their mind slowly lives more by others' reactions than by the values they are pursuing.)}

Khi mức được thích trở thành thước đo chính, chỉ một ánh nhìn lạ cũng đủ làm lòng chao đảo.
\textit{(When being liked becomes the main measure, even one strange look can shake the heart.)}

Giá trị con người không thể đặt trọn trong tay đám đông, vì đám đông đổi rất nhanh.
\textit{(Human worth cannot be placed entirely in the hands of the crowd, because the crowd changes quickly.)}

\end{truyen}

\begin{baihoc}
Người biết mình đang sống vì điều gì sẽ bớt lệ thuộc vào việc mình có đang được thích hay không.
\textit{(Those who know what they live for depend less on whether they are being liked.)}
\end{baihoc}

\begin{ghinhoanh}
\textbf{Short English to remember:}

Know what you live for, and approval loses some of its power.

\end{ghinhoanh}

\begin{tuhocnhanh}
\textbf{Câu hỏi tự học / Self-study questions}

1. Thước đo chính của nhân vật là gì? \textit{What is the character's main measure?}

2. Vì sao thước đo đó bấp bênh? \textit{Why is that measure unstable?}

3. Em muốn đặt giá trị đời mình trên nền nào? \textit{On what foundation do you want to place your worth?}
\end{tuhocnhanh}
